\documentclass[12pt,a4paper]{article}

% ─── PACKAGES ─────────────────────────────────────────────
\usepackage[utf8]{inputenc}
\usepackage[T1]{fontenc}
\usepackage[french]{babel}
\usepackage{amsmath,amssymb,amsthm}
\usepackage{mathtools}
\usepackage{booktabs}
\usepackage{array}
\usepackage{longtable}
\usepackage{graphicx}
\usepackage{xcolor}
\usepackage{hyperref}
\usepackage{enumitem}
\usepackage{geometry}
\usepackage{fancyhdr}
\usepackage{titlesec}
\usepackage{tcolorbox}
\usepackage{caption}
\usepackage{subcaption}
\usepackage{algorithm}
\usepackage{algpseudocode}
\usepackage{tikz}
\usetikzlibrary{arrows.meta, positioning, shapes.geometric, calc}

\geometry{margin=2.5cm}

% ─── COULEURS ─────────────────────────────────────────────
\definecolor{steelblue}{HTML}{4682B4}
\definecolor{darkblue}{HTML}{2C3E50}
\definecolor{lightgray}{HTML}{F5F5F5}
\definecolor{accentgreen}{HTML}{27AE60}
\definecolor{accentorange}{HTML}{E67E22}
\definecolor{accentred}{HTML}{E74C3C}

% ─── STYLE ────────────────────────────────────────────────
\hypersetup{
    colorlinks=true,
    linkcolor=steelblue,
    citecolor=steelblue,
    urlcolor=steelblue,
}

\titleformat{\section}
    {\Large\bfseries\color{darkblue}}{\thesection}{1em}{}
\titleformat{\subsection}
    {\large\bfseries\color{steelblue}}{\thesubsection}{1em}{}
\titleformat{\subsubsection}
    {\normalsize\bfseries\color{steelblue!80}}{\thesubsubsection}{1em}{}

\pagestyle{fancy}
\fancyhf{}
\fancyhead[L]{\small\color{gray}IFRS 9 Risk Cockpit --- M\'ethodologie PD}
\fancyhead[R]{\small\color{gray}\thepage}
\renewcommand{\headrulewidth}{0.4pt}

% ─── ENCADRES ─────────────────────────────────────────────
\newtcolorbox{definitionbox}[1][]{
    colback=steelblue!5,
    colframe=steelblue,
    title=#1,
    fonttitle=\bfseries,
    rounded corners,
    boxrule=0.8pt,
}

\newtcolorbox{formulebox}[1][]{
    colback=lightgray,
    colframe=darkblue,
    title=#1,
    fonttitle=\bfseries,
    rounded corners,
    boxrule=0.8pt,
}

\newtcolorbox{interpretbox}[1][]{
    colback=accentgreen!5,
    colframe=accentgreen!70!black,
    title=#1,
    fonttitle=\bfseries,
    rounded corners,
    boxrule=0.8pt,
}

\newtcolorbox{warningbox}[1][]{
    colback=accentorange!5,
    colframe=accentorange,
    title=#1,
    fonttitle=\bfseries,
    rounded corners,
    boxrule=0.8pt,
}

% ─── DOCUMENT ─────────────────────────────────────────────
\begin{document}

% ─── PAGE DE TITRE ────────────────────────────────────────
\begin{titlepage}
\centering
\vspace*{3cm}

{\Huge\bfseries\color{darkblue} Mod\'elisation de la\\[0.3cm]
Probabilit\'e de D\'efaut (PD)}

\vspace{1cm}

{\LARGE\color{steelblue} M\'ethodologie, Calculs et Interpr\'etations}

\vspace{2cm}

{\large\color{gray}
IFRS 9 Risk Cockpit\\[0.5cm]
Approche multi-mod\`eles : LR\_WoE $\cdot$ TabNet $\cdot$ XGBoost
}

\vspace{3cm}

\begin{tikzpicture}
\node[draw=steelblue, rounded corners=8pt, minimum width=4cm, minimum height=1.2cm, fill=steelblue!10] (lr) {\textbf{LR\_WoE}};
\node[draw=steelblue, rounded corners=8pt, minimum width=4cm, minimum height=1.2cm, fill=steelblue!10, right=1.5cm of lr] (tn) {\textbf{TabNet}};
\node[draw=steelblue, rounded corners=8pt, minimum width=4cm, minimum height=1.2cm, fill=steelblue!10, right=1.5cm of tn] (xgb) {\textbf{XGBoost}};

\node[below=0.3cm of lr, text=gray, font=\small] {Transparence};
\node[below=0.3cm of tn, text=gray, font=\small] {Explicabilit\'e};
\node[below=0.3cm of xgb, text=gray, font=\small] {Performance};
\end{tikzpicture}

\vfill

{\large\today}

\end{titlepage}

\tableofcontents
\newpage

% ══════════════════════════════════════════════════════════
\section{Introduction et Contexte R\'eglementaire}
% ══════════════════════════════════════════════════════════

\subsection{Pourquoi mod\'eliser la PD~?}

La norme IFRS~9 impose aux \'etablissements financiers de provisionner les pertes de cr\'edit attendues (\textit{Expected Credit Loss}, ECL) d\`es l'origination du pr\^et. La formule fondamentale est~:

\begin{formulebox}[Formule ECL fondamentale]
\begin{equation}
    \text{ECL} = \text{PD} \times \text{LGD} \times \text{EAD} \times \text{DF}
    \label{eq:ecl}
\end{equation}
o\`u~:
\begin{itemize}[noitemsep]
    \item $\text{PD}$ = Probabilit\'e de D\'efaut (notre sujet)
    \item $\text{LGD}$ = \textit{Loss Given Default} (perte en cas de d\'efaut)
    \item $\text{EAD}$ = \textit{Exposure At Default} (exposition au moment du d\'efaut)
    \item $\text{DF}$ = Facteur d'actualisation
\end{itemize}
\end{formulebox}

La PD est le param\`etre le plus sensible de cette \'equation~: une erreur de 1 point de pourcentage sur la PD se propage directement dans l'ECL, et donc dans le compte de r\'esultat de la banque.

\subsection{Exigences r\'eglementaires (BCE / EBA)}

Le r\'egulateur europ\'een impose plusieurs exigences sur la mod\'elisation PD~:

\begin{enumerate}[noitemsep]
    \item \textbf{Model Risk Management}~: comparer plusieurs familles de mod\`eles pour \'eviter le risque de mod\`ele unique
    \item \textbf{Interpr\'etabilit\'e}~: pouvoir expliquer les d\'ecisions de cr\'edit \`a un client ou un auditeur
    \item \textbf{Calibration}~: les PD pr\'edites doivent correspondre aux taux de d\'efaut observ\'es
    \item \textbf{Stabilit\'e}~: le mod\`ele doit \^etre robuste dans le temps (monitoring PSI/CSI)
\end{enumerate}

\subsection{Choix de trois familles de mod\`eles}

Pour r\'epondre \`a ces exigences, nous d\'eployons trois mod\`eles compl\'ementaires~:

\begin{center}
\begin{tabular}{lllp{5cm}}
\toprule
\textbf{Mod\`ele} & \textbf{Famille} & \textbf{R\^ole} & \textbf{Public cible} \\
\midrule
LR\_WoE & Lin\'eaire & Transparence, audit & R\'egulateurs, auditeurs \\
TabNet & Deep Learning & Explicabilit\'e individuelle & Comit\'e de cr\'edit, CRO \\
XGBoost & Ensemble d'arbres & Performance maximale & \'Equipe risque, production \\
\bottomrule
\end{tabular}
\end{center}

\begin{interpretbox}[Interpr\'etation]
Chaque mod\`ele r\'epond \`a un besoin diff\'erent~: un mod\`ele qu'on \textbf{comprend} (LR\_WoE), un qu'on \textbf{utilise} en production (XGBoost), un qu'on \textbf{explique} individuellement (TabNet). C'est le trio classique en banque, conforme aux bonnes pratiques BCE/EBA.
\end{interpretbox}


% ══════════════════════════════════════════════════════════
\section{Donn\'ees et G\'en\'erateur Synth\'etique (DGP v4.5)}
% ══════════════════════════════════════════════════════════

\subsection{Architecture du DGP}

Les donn\'ees sont produites par un \textit{Data Generating Process} (DGP) synth\'etique con\c{c}u pour reproduire les caract\'eristiques statistiques d'un portefeuille de cr\'edit corporate. L'architecture suit une d\'ecomposition 50/30/20~:

\begin{center}
\begin{tikzpicture}[
    box/.style={draw=steelblue, rounded corners=5pt, minimum width=3.5cm, minimum height=1cm, fill=steelblue!10, font=\small},
    arrow/.style={-Stealth, thick, steelblue}
]
\node[box] (mono) at (0,0) {50\% Monotones\\par morceaux};
\node[box] (cond) at (5.5,0) {30\% Amplifications\\conditionnelles};
\node[box] (conj) at (11,0) {20\% Conjunctives\\(interactions)};

\node[box, fill=darkblue!10, draw=darkblue, minimum width=12cm] (logit) at (5.5,-2) {$\text{logit}(\text{PD}) = \beta_0 + f_{\text{mono}}(\mathbf{x}) + f_{\text{cond}}(\mathbf{x}) + f_{\text{conj}}(\mathbf{x}) + \varepsilon$};

\draw[arrow] (mono) -- (logit);
\draw[arrow] (cond) -- (logit);
\draw[arrow] (conj) -- (logit);
\end{tikzpicture}
\end{center}

\begin{itemize}
    \item \textbf{Monotones par morceaux (50\%)}~: effets lin\'eaires par segments (ex.~: \texttt{debt\_ratio} augmente la PD lin\'eairement par tranches)
    \item \textbf{Amplifications conditionnelles (30\%)}~: effets non-lin\'eaires activ\'es sous condition (ex.~: seuil de dette \`a 60\% avec effet quadratique au-del\`a)
    \item \textbf{Conjunctives (20\%)}~: interactions entre variables (ex.~: \texttt{debt\_ratio} $\times$ \texttt{interest\_rate})
\end{itemize}

\subsection{Variables g\'en\'er\'ees}

\begin{center}
\begin{tabular}{llp{7cm}}
\toprule
\textbf{Variable} & \textbf{Type} & \textbf{Description} \\
\midrule
\texttt{revenue} & Num\'erique & Chiffre d'affaires (log-normal par secteur) \\
\texttt{ebitda} & Num\'erique & R\'esultat op\'erationnel (corr\'el\'e au revenue) \\
\texttt{debt\_ratio} & Num\'erique & Ratio d'endettement (distribution Gamma) \\
\texttt{credit\_score} & Num\'erique & Score de cr\'edit [300, 850] \\
\texttt{dpd} & Num\'erique & Jours de retard de paiement \\
\texttt{utilization\_rate} & Num\'erique & Taux d'utilisation de la ligne \\
\texttt{loan\_amount} & Num\'erique & Montant du pr\^et \\
\texttt{collateral} & Num\'erique & Valeur de la garantie \\
\texttt{sector} & Cat\'egorielle & Secteur \'economique (5 secteurs) \\
\texttt{loan\_type} & Cat\'egorielle & Type de pr\^et (Revolving / Term) \\
\texttt{loan\_to\_revenue} & Engin. & $= \texttt{loan\_amount} / \texttt{revenue}$ \\
\texttt{collateral\_coverage} & Engin. & $= \texttt{collateral} / \texttt{loan\_amount}$ \\
\bottomrule
\end{tabular}
\end{center}

\subsection{S\'eparation entra\^inement / portfolio}

\begin{warningbox}[Gestion des seeds]
\begin{itemize}[noitemsep]
    \item \textbf{Entra\^inement}~: 1\,500\,000 lignes, seed = 42
    \item \textbf{Portfolio (dashboard)}~: 30\,000 lignes, seed = 123
    \item Les deux seeds sont \textbf{toujours diff\'erentes} pour \'eviter la circularit\'e (le mod\`ele ne doit jamais \^etre \'evalu\'e sur ses donn\'ees d'entra\^inement)
\end{itemize}
\end{warningbox}


% ══════════════════════════════════════════════════════════
\section{Mod\`ele 1~: R\'egression Logistique avec WoE (LR\_WoE)}
\label{sec:lr_woe}
% ══════════════════════════════════════════════════════════

Ce mod\`ele est la \textbf{scorecard bancaire classique}. C'est la r\'ef\'erence de l'industrie depuis les ann\'ees 1990 (Siddiqi, 2006). Son avantage principal~: \textbf{transparence totale}. Chaque d\'ecision peut \^etre expliqu\'ee par une formule lin\'eaire.

\subsection{\'Etape 1~: Weight of Evidence (WoE)}

\subsubsection{Principe}

Le WoE transforme chaque variable num\'erique en une mesure de \textbf{pouvoir discriminant} par tranche. L'id\'ee~: d\'ecouper la variable en bins (intervalles), puis mesurer si chaque bin contient proportionnellement plus de ``bons'' ou de ``mauvais'' clients.

\begin{definitionbox}[D\'efinition -- Weight of Evidence]
Pour un bin $i$ d'une variable $X$~:
\begin{equation}
    \text{WoE}_i = \ln\left(\frac{\%\;\text{non-d\'efauts dans le bin}_i}{\%\;\text{d\'efauts dans le bin}_i}\right)
    = \ln\left(\frac{n_{i}^{(\text{good})} / N^{(\text{good})}}{n_{i}^{(\text{bad})} / N^{(\text{bad})}}\right)
    \label{eq:woe}
\end{equation}
\end{definitionbox}

\textbf{Interpr\'etation intuitive}~:
\begin{itemize}[noitemsep]
    \item $\text{WoE} > 0$~: le bin contient proportionnellement plus de bons clients $\Rightarrow$ tranche s\^ure
    \item $\text{WoE} < 0$~: le bin contient proportionnellement plus de mauvais clients $\Rightarrow$ tranche risqu\'ee
    \item $\text{WoE} = 0$~: le bin a la m\^eme proportion que la population globale $\Rightarrow$ neutre
\end{itemize}

\subsubsection{Exemple concret~: \texttt{debt\_ratio}}

\begin{center}
\begin{tabular}{ccccc}
\toprule
\textbf{Bin} & \textbf{\% Bons} & \textbf{\% Mauvais} & \textbf{WoE} & \textbf{Interpr\'etation} \\
\midrule
$[0.0,\; 0.2]$ & 25\% & 8\% & $+1.14$ & Tr\`es s\^ur \\
$[0.2,\; 0.4]$ & 22\% & 15\% & $+0.38$ & S\^ur \\
$[0.4,\; 0.6]$ & 20\% & 22\% & $-0.10$ & L\'eg\`erement risqu\'e \\
$[0.6,\; 0.8]$ & 18\% & 25\% & $-0.33$ & Risqu\'e \\
$[0.8,\; 1.0]$ & 15\% & 30\% & $-0.69$ & Tr\`es risqu\'e \\
\bottomrule
\end{tabular}
\end{center}

\subsubsection{Lissage de Laplace ($\varepsilon$-smoothing)}

Un bin peut contenir z\'ero d\'efaut (ou z\'ero non-d\'efaut), ce qui donnerait $\ln(x/0) = +\infty$. Pour \'eviter ce probl\`eme, on applique un \textbf{lissage de Laplace}~:

\begin{formulebox}[WoE avec lissage]
\begin{equation}
    \text{WoE}_i = \ln\left(
        \frac{(n_{i}^{(\text{good})} + \varepsilon)\;/\;(N^{(\text{good})} + 2\varepsilon)}
             {(n_{i}^{(\text{bad})} + \varepsilon)\;/\;(N^{(\text{bad})} + 2\varepsilon)}
    \right)
    \label{eq:woe_smooth}
\end{equation}
Avec $\varepsilon = 0.5$ (valeur par d\'efaut, configurable).
\end{formulebox}

Le terme $2\varepsilon$ au d\'enominateur assure que les proportions somment toujours \`a 1 apr\`es lissage. C'est la correction de Jeffreys, standard en statistique bay\'esienne.

\subsection{\'Etape 2~: Monotonicit\'e (PAV)}

\subsubsection{Probl\`eme}

Les WoE bruts peuvent \^etre \textbf{non-monotones} par hasard statistique. Par exemple, le bin $[0.4, 0.6]$ pourrait avoir un taux de d\'efaut \textit{inf\'erieur} au bin $[0.2, 0.4]$, simplement par fluctuation d'\'echantillonnage. Or, pour une scorecard r\'eglementaire, on exige que~:

\begin{quote}
\textit{``Plus la dette augmente, plus le risque augmente''} $\Rightarrow$ WoE strictement d\'ecroissant.
\end{quote}

\subsubsection{Solution~: Pool Adjacent Violators (PAV)}

L'algorithme PAV (aussi appel\'e \textit{Isotonic Regression}) corrige les violations de monotonicit\'e en \textbf{fusionnant les bins adjacents} qui violent l'ordre attendu~:

\begin{algorithm}[H]
\caption{Pool Adjacent Violators (PAV)}
\begin{algorithmic}[1]
\Require Taux de d\'efaut par bin $\{r_1, r_2, \ldots, r_k\}$, direction attendue
\Ensure Taux monotones $\{r_1', r_2', \ldots, r_m'\}$ avec $m \leq k$
\State D\'eterminer la direction via corr\'elation de Spearman
\State Appliquer la r\'egression isotonique (sklearn)
\For{chaque paire de bins adjacents avec le m\^eme taux corrig\'e}
    \State Fusionner les deux bins (somme des effectifs)
    \State Supprimer le breakpoint interm\'ediaire
\EndFor
\State \Return Bins fusionn\'es (monotones par construction)
\end{algorithmic}
\end{algorithm}

\subsubsection{D\'etection automatique de la direction}

La direction (croissante ou d\'ecroissante) est d\'etect\'ee par la \textbf{corr\'elation de Spearman} entre la m\'ediane du bin et le taux de d\'efaut~:

\begin{itemize}[noitemsep]
    \item $\rho_{\text{Spearman}} < 0$~: quand la variable augmente, le taux de d\'efaut diminue $\Rightarrow$ WoE croissant (variable protectrice)
    \item $\rho_{\text{Spearman}} > 0$~: quand la variable augmente, le taux de d\'efaut augmente $\Rightarrow$ WoE d\'ecroissant (variable de risque)
\end{itemize}

\subsection{\'Etape 3~: Nombre minimum de d\'efauts par bin}

Un bin avec seulement 3 d\'efauts sur 10\,000 observations n'est pas statistiquement fiable. On impose un \textbf{minimum de d\'efauts par bin} (param\`etre \texttt{min\_events\_per\_bin} $= 20$).

Si un bin contient moins de 20 d\'efauts, il est \textbf{fusionn\'e avec le bin voisin} le plus proche (celui ayant le moins de d\'efauts). Ce processus est it\'eratif jusqu'\`a ce que tous les bins respectent le seuil.

\begin{interpretbox}[Pourquoi 20~?]
Avec 20 d\'efauts par bin, l'erreur standard du taux de d\'efaut est~:
$$\text{SE} = \sqrt{\frac{p(1-p)}{n}} \approx \sqrt{\frac{0.05 \times 0.95}{20}} \approx 4.9\%$$
C'est un compromis entre pr\'ecision statistique et granularit\'e des bins.
\end{interpretbox}

\subsection{\'Etape 3b~: Granularit\'e minimale post-PAV}

L'algorithme PAV peut cr\'eer des bins minuscules en fusionnant des violations de monotonicit\'e. Un bin repr\'esentant moins de 5\% de la population totale n'est pas statistiquement repr\'esentatif et risque de provoquer du sur-apprentissage.

\begin{warningbox}[Contrainte de granularit\'e]
Apr\`es PAV, tout bin contenant \textbf{moins de 5\% de la population totale} (param\`etre \texttt{min\_bin\_pct} $= 0.05$) est fusionn\'e avec le bin voisin dont le \textbf{WoE est le plus proche} (et non le voisin num\'erique imm\'ediat). Cette technique pr\'eserve la s\'eparation optimale du mod\`ele en regroupant les segments les plus similaires en termes de risque.
\end{warningbox}

\subsection{\'Etape 3c~: Traitement du bin NaN (WoE imputation)}

Les valeurs manquantes (NaN) sont isol\'ees dans un \textbf{bin explicite} avec leur propre WoE calcul\'e ind\'ependamment. Cela \'evite l'imputation par la m\'ediane, qui effacerait un signal potentiellement informatif (ex.~: un champ vide peut indiquer une anomalie dans le processus d'octroi).

Si le bin NaN contient \textbf{moins de 5\% de la population}, il est fusionn\'e avec le bin num\'erique dont le WoE est le plus proche~:

\begin{quote}
\textit{Fusion par proximit\'e WoE (``WoE imputation'', Anderson 2007)~: on choisit le bin de destination par $\min_b |WoE_{\text{NaN}} - WoE_b|$, et non par proximit\'e num\'erique. Cela garantit que les NaN sont class\'es dans la cat\'egorie de risque la plus similaire.}
\end{quote}

\subsection{\'Etape 4~: Information Value (IV) --- S\'election de variables}

L'IV mesure le \textbf{pouvoir discriminant total} d'une variable sur l'ensemble de ses bins~:

\begin{formulebox}[Information Value]
\begin{equation}
    \text{IV} = \sum_{i=1}^{k} \left(\%\;\text{bons}_i - \%\;\text{mauvais}_i\right) \times \text{WoE}_i
    \label{eq:iv}
\end{equation}
\end{formulebox}

\begin{center}
\begin{tabular}{cl}
\toprule
\textbf{IV} & \textbf{Interpr\'etation} \\
\midrule
$< 0.02$ & Pas de pouvoir pr\'edictif $\Rightarrow$ \textbf{\'elimin\'ee} \\
$0.02 - 0.10$ & Faible \\
$0.10 - 0.30$ & Moyen \\
$0.30 - 0.50$ & Fort \\
$> 0.50$ & Suspect (possible fuite de donn\'ees) \\
\bottomrule
\end{tabular}
\end{center}

Les variables avec $\text{IV} < 0.02$ sont automatiquement exclues du mod\`ele.

\begin{interpretbox}[Interpr\'etation de l'IV]
L'IV est en fait la \textbf{divergence de Kullback-Leibler sym\'etris\'ee} entre la distribution des bons et des mauvais clients. Plus cette divergence est grande, plus la variable s\'epare bien les deux populations.
\end{interpretbox}

\subsection{\'Etape 5~: D\'efinition de la cible}

\begin{warningbox}[Convention de la variable cible]
Le mod\`ele estime explicitement $P(\text{D\'efaut} = 1 \mid \mathbf{X})$, c'est-\`a-dire la probabilit\'e de d\'efaut. Certains syst\`emes (notamment des scorecards am\'ericaines) mod\'elisent $P(\text{Non-D\'efaut} = 1)$, ce qui \textbf{inverse tous les signes}. Ici, la cible est \texttt{default\_flag} $\in \{0, 1\}$ o\`u $1 =$ d\'efaut. Toutes les formules, conventions de signe ($\beta < 0$) et directions WoE sont coh\'erentes avec cette convention.
\end{warningbox}

\subsection{\'Etape 5b~: R\'egression Logistique}

Apr\`es l'encodage WoE, chaque variable num\'erique est remplac\'ee par un scalaire (le WoE de sa tranche). Le mod\`ele est alors une simple r\'egression logistique~:

\begin{formulebox}[Mod\`ele logistique]
\begin{equation}
    \text{logit}(\text{PD}) = \ln\left(\frac{\text{PD}}{1 - \text{PD}}\right) = \beta_0 + \sum_{j=1}^{p} \beta_j \cdot \text{WoE}_j(x_j)
    \label{eq:logit}
\end{equation}
\begin{equation}
    \text{PD} = \sigma\left(\beta_0 + \sum_{j=1}^{p} \beta_j \cdot \text{WoE}_j(x_j)\right) = \frac{1}{1 + e^{-(\beta_0 + \sum \beta_j \cdot \text{WoE}_j)}}
    \label{eq:pd}
\end{equation}
\end{formulebox}

\subsubsection{Filtrage de la multicolin\'earit\'e (VIF)}

Avant de v\'erifier les signes des $\beta$, on filtre les variables pr\'esentant une \textbf{multicolin\'earit\'e s\'ev\`ere}. Un coefficient $\beta_j > 0$ inattendu (alors que le WoE est monotone) est souvent le sympt\^ome d'un \textit{effet de masque} entre variables corr\'el\'ees, plut\^ot que d'une mauvaise qualit\'e intrins\`eque de la variable.

\begin{formulebox}[Variance Inflation Factor]
\begin{equation}
    \text{VIF}_j = \frac{1}{1 - R_j^2}
    \label{eq:vif}
\end{equation}
o\`u $R_j^2$ est le coefficient de d\'etermination de la r\'egression de $\text{WoE}_j$ sur les autres $\text{WoE}_{k \neq j}$.
\end{formulebox}

\begin{center}
\begin{tabular}{cl}
\toprule
\textbf{VIF} & \textbf{Interpr\'etation} \\
\midrule
$1$ & Pas de corr\'elation \\
$1 - 5$ & Corr\'elation mod\'er\'ee (acceptable) \\
$> 5$ & Multicolin\'earit\'e forte $\Rightarrow$ \textbf{variable supprim\'ee} \\
\bottomrule
\end{tabular}
\end{center}

Le processus est it\'eratif~: on supprime la variable avec le VIF le plus \'elev\'e ($> 5$), on recalcule les VIF, et on r\'ep\`ete jusqu'\`a ce que tous les VIF soient acceptables. Cela stabilise les signes des $\beta$ sans rejeter arbitrairement des variables pr\'edictives.

\subsubsection{Contrainte $\beta \leq 0$}

Chaque coefficient $\beta_j$ doit \^etre \textbf{n\'egatif ou nul}. La justification est la suivante~:

\begin{itemize}
    \item Un WoE positif signifie ``tranche s\^ure'' $\Rightarrow$ doit \textbf{diminuer} la PD $\Rightarrow$ $\beta < 0$
    \item Un WoE n\'egatif signifie ``tranche risqu\'ee'' $\Rightarrow$ $\beta < 0 \times \text{WoE} < 0 = \text{positif}$ $\Rightarrow$ \textbf{augmente} la PD \checkmark
\end{itemize}

Apr\`es le filtrage VIF, si un $\beta_j > 0$ subsiste (rare une fois la multicolin\'earit\'e trait\'ee), la variable est \textbf{supprim\'ee} et le mod\`ele est r\'e-estim\'e. Le processus est it\'eratif jusqu'\`a ce que tous les $\beta_j \leq 0$.

\begin{warningbox}[Alternative~: optimisation sous contrainte]
Une approche plus \'el\'egante serait d'estimer directement la LR sous contrainte de signe via \texttt{scipy.optimize.minimize} avec \textit{bounds}. L'\'elimination it\'erative apr\`es filtrage VIF est cependant standard en pratique bancaire (Anderson, 2007) et donne des r\'esultats \'equivalents dans la majorit\'e des cas.
\end{warningbox}

\subsection{\'Etape 6~: Correction du Prior}

Le mod\`ele est entra\^in\'e avec \texttt{class\_weight="balanced"}, ce qui attribue des poids inversement proportionnels \`a la fr\'equence de chaque classe. Cela revient \`a entra\^iner sur un \'echantillon 50/50, ce qui am\'eliore l'apprentissage des patterns (les d\'efauts repr\'esentent typiquement 5\% du portefeuille).

Cependant, les probabilit\'es pr\'edites sont alors \textbf{biais\'ees vers 50\%}. On corrige l'intercept~:

\begin{formulebox}[Correction du prior (King \& Zeng, 2001)]
\begin{equation}
    \beta_{0,\text{corrig\'e}} = \beta_0 + \underbrace{\ln\left(\frac{\pi_{\text{pop}}}{1 - \pi_{\text{pop}}}\right)}_{\text{logit pr\'evalence r\'eelle}} - \underbrace{\ln\left(\frac{\pi_{\text{train}}}{1 - \pi_{\text{train}}}\right)}_{\text{logit pr\'evalence entra\^inement}}
    \label{eq:prior}
\end{equation}
Avec $\pi_{\text{pop}} = 0.05$ (taux de d\'efaut r\'eel) et $\pi_{\text{train}} = 0.5$ (\'echantillonnage balanc\'e).
\end{formulebox}

\begin{interpretbox}[Pourquoi cette formule fonctionne~?]
La r\'egression logistique avec sur-\'echantillonnage ne modifie que l'intercept, pas les pentes (propri\'et\'e connue du maximum de vraisemblance conditionnelle). La correction revient \`a ``d\'ecaler'' toutes les PD pr\'edites du bon montant pour retrouver le niveau moyen r\'eel. En terme num\'erique~:
$$\text{correction} = \ln\left(\frac{0.05}{0.95}\right) - \underbrace{\ln\left(\frac{0.5}{0.5}\right)}_{= \ln(1) = 0} = -2.944 - 0 = -2.944$$
Toutes les PD sont donc d\'ecal\'ees vers le bas d'environ 3 unit\'es logit.

\textbf{Condition d'annulation}~: le second terme ne s'annule que si $\pi_{\text{train}} = 0.5$ exactement. C'est le cas avec \texttt{class\_weight="balanced"} (les poids \'equilibrent parfaitement les classes). Si une autre strat\'egie de r\'e-\'echantillonnage est utilis\'ee (SMOTE, undersampling \`a un ratio $\neq$ 50/50), il faut mettre \`a jour $\pi_{\text{train}}$ en cons\'equence, et la formule g\'en\'erale s'applique pleinement.
\end{interpretbox}

\subsection{\'Etape 7~: Conversion en Scorecard}

Pour les op\'erationnels, une PD de 3.2\% est moins intuitive qu'un \textbf{score de 720 points}. On convertit donc la PD en score entier via une transformation log-odds~:

\begin{formulebox}[Scorecard (Siddiqi, 2006)]
\begin{equation}
    \text{Score} = \underbrace{\text{Offset}}_{\text{point d'ancrage}} + \underbrace{\text{Factor}}_{\text{\'echelle}} \times \ln\left(\frac{\text{PD}}{1 - \text{PD}}\right)
    \label{eq:score}
\end{equation}
Avec~:
\begin{align}
    \text{Factor} &= \frac{-\text{PDO}}{\ln(2)} = \frac{-20}{\ln(2)} \approx -28.85 \label{eq:factor} \\
    \text{Offset} &= \text{Target\_Score} - \text{Factor} \times \ln(\text{Target\_Odds}) \label{eq:offset}
\end{align}
\end{formulebox}

\textbf{Param\`etres de calibration}~:
\begin{itemize}[noitemsep]
    \item \textbf{Target\_Score} $= 600$~: score d'un client au odds de r\'ef\'erence
    \item \textbf{Target\_Odds} $= 50$~: ratio de cotes de r\'ef\'erence (50 bons pour 1 mauvais, soit PD $\approx 2\%$)
    \item \textbf{PDO} $= 20$~: \textit{Points to Double the Odds} --- chaque $+20$ points de score signifie que les odds sont multipli\'es par 2
\end{itemize}

\begin{center}
\begin{tabular}{cccc}
\toprule
\textbf{PD} & \textbf{Odds (bons:mauvais)} & \textbf{Score} & \textbf{Cat\'egorie} \\
\midrule
1\% & 99:1 & 846 & Excellent \\
2\% & 49:1 & 713 & Bon \\
5\% & 19:1 & 628 & Moyen \\
10\% & 9:1 & 576 & Risqu\'e \\
20\% & 4:1 & 513 & Tr\`es risqu\'e \\
\bottomrule
\end{tabular}
\end{center}

\subsection{Pipeline complet LR\_WoE}

\begin{center}
\begin{tikzpicture}[
    box/.style={draw=steelblue, rounded corners=4pt, minimum width=3.2cm, minimum height=0.9cm, fill=steelblue!8, font=\small, text=darkblue, align=center},
    newbox/.style={draw=accentgreen!70!black, rounded corners=4pt, minimum width=3.2cm, minimum height=0.9cm, fill=accentgreen!8, font=\small, text=darkblue, align=center},
    arrow/.style={-Stealth, thick, steelblue!70}
]
\node[box] (bin) at (0,0) {D\'ecoupage\\en bins};
\node[box] (woe) at (4,0) {Calcul WoE\\+ lissage $\varepsilon$};
\node[box] (pav) at (8,0) {PAV\\(monotonicit\'e)};
\node[newbox] (minbin) at (12,0) {Granularit\'e\\$\geq 5\%$};
\node[box] (iv) at (0,-2) {S\'election IV\\$\geq 0.02$};
\node[newbox] (vif) at (4,-2) {Filtrage VIF\\$\leq 5$};
\node[box] (lr) at (8,-2) {R\'egression\\logistique};
\node[box] (beta) at (12,-2) {\'Elimination\\$\beta > 0$};
\node[box] (prior) at (3,-4) {Correction\\du prior};
\node[box] (score) at (9,-4) {Scorecard\\(points)};

\draw[arrow] (bin) -- (woe);
\draw[arrow] (woe) -- (pav);
\draw[arrow] (pav) -- (minbin);
\draw[arrow] (minbin) -- ++(0,-0.5) -| (iv);
\draw[arrow] (iv) -- (vif);
\draw[arrow] (vif) -- (lr);
\draw[arrow] (lr) -- (beta);
\draw[arrow] (beta) -- ++(0,-0.5) -| (prior);
\draw[arrow] (prior) -- (score);
\end{tikzpicture}
\end{center}

\begin{interpretbox}[Nouvelles \'etapes (en vert)]
Deux \'etapes ont \'et\'e ajout\'ees pour renforcer la robustesse r\'eglementaire~:
\begin{itemize}[noitemsep]
    \item \textbf{Granularit\'e $\geq 5\%$}~: fusion post-PAV des micro-bins par proximit\'e WoE (\'evite le sur-apprentissage)
    \item \textbf{Filtrage VIF $\leq 5$}~: \'elimination it\'erative des variables multicolin\'eaires \textit{avant} la contrainte de signe (stabilise les $\beta$)
\end{itemize}
\end{interpretbox}


% ══════════════════════════════════════════════════════════
\section{Mod\`ele 2~: TabNet (Deep Learning Tabulaire)}
\label{sec:tabnet}
% ══════════════════════════════════════════════════════════

\subsection{Pourquoi du deep learning sur des donn\'ees tabulaires~?}

Le deep learning domine sur les images (CNN), le texte (Transformers) et l'audio (WaveNet). Sur les \textbf{donn\'ees tabulaires}, la question est l\'egitime~: les m\'ethodes ensemblistes (XGBoost, Random Forest) sont souvent aussi bonnes, voire meilleures.

TabNet (Arik \& Pfister, 2019, Google Research) est sp\'ecifiquement con\c{c}u pour les donn\'ees tabulaires. Son avantage unique~: des \textbf{masques d'attention} qui montrent, \textit{pour chaque client individuellement}, quelles variables ont compt\'e dans la d\'ecision.

\subsection{Architecture}

\begin{definitionbox}[Architecture TabNet]
TabNet proc\`ede en $N_{\text{steps}}$ \'etapes s\'equentielles. \`A chaque \'etape~:
\begin{enumerate}[noitemsep]
    \item Un \textbf{masque d'attention} s\'electionne un sous-ensemble de variables
    \item Les variables s\'electionn\'ees passent dans un r\'eseau dense (\textit{Feature Transformer})
    \item La sortie est accumul\'ee pour la pr\'ediction finale
\end{enumerate}
\end{definitionbox}

\subsubsection{Masques d'attention sparses}

\`A chaque \'etape $s$, TabNet calcule un masque d'attention $\mathbf{M}^{(s)} \in [0, 1]^D$ (o\`u $D$ est le nombre de features)~:

\begin{equation}
    \mathbf{M}^{(s)} = \text{Sparsemax}\left(\mathbf{P}^{(s-1)} \cdot h_s\left(\mathbf{a}^{(s-1)}\right)\right)
    \label{eq:attention}
\end{equation}

o\`u~:
\begin{itemize}[noitemsep]
    \item $\mathbf{P}^{(s-1)}$~: prior scales (force l'exploration de nouvelles features)
    \item $h_s$~: r\'eseau appris (fully connected + BN)
    \item $\text{Sparsemax}$~: version sparse de softmax (pousse les poids faibles \`a exactement z\'ero)
\end{itemize}

Le coefficient $\gamma$ (= 1.3 dans notre config) contr\^ole la r\'eutilisation des features~: $\gamma = 1$ interdit la r\'eutilisation, $\gamma > 1$ l'autorise partiellement.

\subsubsection{Feature Transformer}

Chaque \'etape transforme les features s\'electionn\'ees via~:

\begin{equation}
    \mathbf{h}^{(s)} = \text{ReLU}\left(\text{BN}\left(\text{FC}_{n_d}\left(\mathbf{M}^{(s)} \odot \mathbf{x}\right)\right)\right)
    \label{eq:ft}
\end{equation}

Avec $\odot$ le produit \'el\'ement par \'el\'ement (masquage), $n_d = 64$ la dimension de sortie.

\subsubsection{Pr\'ediction finale}

La sortie est la somme des contributions de chaque \'etape~:

\begin{equation}
    \hat{y} = \sigma\left(\sum_{s=1}^{N_{\text{steps}}} \text{ReLU}\left(\mathbf{h}^{(s)}\right) \cdot \mathbf{W}_{\text{final}}\right)
    \label{eq:tabnet_pred}
\end{equation}

\subsection{Hyperparamètres}

\begin{center}
\begin{tabular}{lcp{7cm}}
\toprule
\textbf{Param\`etre} & \textbf{Valeur} & \textbf{R\^ole} \\
\midrule
$n_d$ (decision) & 64 & Dimension des couches de d\'ecision \\
$n_a$ (attention) & 64 & Dimension des couches d'attention \\
$N_{\text{steps}}$ & 5 & Nombre d'\'etapes d'attention s\'equentielles \\
$\gamma$ & 1.3 & Coefficient de r\'eutilisation des features \\
$\lambda_{\text{sparse}}$ & $10^{-3}$ & P\'enalit\'e de sparsit\'e (encourage les masques \`a z\'ero) \\
Batch size & 8\,192 & Taille des mini-lots \\
Virtual batch & 512 & Ghost Batch Normalization \\
\bottomrule
\end{tabular}
\end{center}

Nombre total de param\`etres~: $\sim$120\,000 (comparable \`a un petit r\'eseau de neurones).

\subsection{Entra\^inement et Early Stopping}

\subsubsection{Qu'est-ce qu'une \'epoque~?}

\begin{definitionbox}[\'Epoque]
Une \textbf{\'epoque} = une passe compl\`ete sur l'ensemble des donn\'ees d'entra\^inement. Avec 1\,500\,000 lignes et un batch de 8\,192~:
$$1\;\text{\'epoque} = \left\lceil\frac{1\,500\,000}{8\,192}\right\rceil = 184\;\text{batches (mini-ajustements des poids)}$$
\end{definitionbox}

\subsubsection{Early stopping}

On r\'eserve 20\% des donn\'ees comme \textbf{validation}. Apr\`es chaque \'epoque, on mesure l'AUC sur ce jeu de validation~:

\begin{itemize}[noitemsep]
    \item Si l'AUC s'am\'eliore $\Rightarrow$ on sauvegarde les poids (``meilleur mod\`ele'')
    \item Si l'AUC ne s'am\'eliore plus pendant \textbf{patience = 20} \'epoques cons\'ecutives $\Rightarrow$ \textbf{arr\^et}
    \item Le mod\`ele final utilise les poids de la \textbf{meilleure \'epoque}, pas de la derni\`ere
\end{itemize}

\begin{interpretbox}[R\'esultat observ\'e]
Sur 1.5M lignes~: arr\^et \`a l'\'epoque 40, meilleur score \`a l'\'epoque 20 ($\text{AUC} = 0.814$). Le mod\`ele a donc vu chaque client environ 20 fois avant d'atteindre son optimum.
\end{interpretbox}

\subsubsection{Co\^ut computationnel}

\begin{center}
\begin{tabular}{lrr}
\toprule
\textbf{Op\'eration} & \textbf{Calculs} & \textbf{Cumul} \\
\midrule
1 forward + backward (1 client) & $\sim$720\,000 FLOPs & -- \\
1 batch (8\,192 clients) & $\sim$5.9 milliards & -- \\
1 \'epoque (1.5M clients) & $\sim$1 TFLOP & 1 TFLOP \\
20 \'epoques utiles & -- & \textbf{$\sim$20 TFLOP} \\
\bottomrule
\end{tabular}
\end{center}

\begin{center}
\begin{tabular}{lcc}
\toprule
\textbf{Hardware} & \textbf{Puissance} & \textbf{Temps r\'eel} \\
\midrule
CPU (Intel i9) & $\sim$0.5 TFLOPS & $\sim$10--15 min \\
GPU (RTX 5070 Ti) & $\sim$35 TFLOPS & $\sim$2--3 min \\
\bottomrule
\end{tabular}
\end{center}

\subsection{Feature importance individuelle}

L'avantage unique de TabNet~: les masques d'attention donnent une importance \textbf{par client}~:

$$\text{Importance}_j^{(\text{client}\;i)} = \sum_{s=1}^{N_{\text{steps}}} \eta^{(s)} \cdot M_{j}^{(s)}(i)$$

o\`u $\eta^{(s)}$ est le coefficient de l'\'etape $s$ et $M_{j}^{(s)}(i)$ le masque d'attention de la feature $j$ pour le client $i$ \`a l'\'etape $s$.

\begin{interpretbox}[Utilisation m\'etier]
Devant un comit\'e de cr\'edit~: ``\textit{Pour le client Durand SA, la d\'ecision repose \`a 45\% sur son ratio d'endettement et \`a 30\% sur son retard de paiement. Les autres variables sont secondaires.}'' --- Ce type d'explication individuelle est impossible avec XGBoost ou la r\'egression logistique.
\end{interpretbox}


% ══════════════════════════════════════════════════════════
\section{Mod\`ele 3~: XGBoost (Gradient Boosting)}
\label{sec:xgboost}
% ══════════════════════════════════════════════════════════

\subsection{Principe du Gradient Boosting}

XGBoost (\textit{eXtreme Gradient Boosting}, Chen \& Guestrin, 2016) est un ensemble de \textbf{petits arbres de d\'ecision} entra\^in\'es s\'equentiellement. Chaque nouvel arbre corrige les erreurs des arbres pr\'ec\'edents.

\begin{formulebox}[Gradient Boosting]
\begin{equation}
    \hat{y}_i = \sum_{t=1}^{T} \eta \cdot f_t(\mathbf{x}_i)
    \label{eq:xgb}
\end{equation}
o\`u~:
\begin{itemize}[noitemsep]
    \item $T = 200$~: nombre d'arbres
    \item $\eta = 0.05$~: learning rate (pas d'apprentissage)
    \item $f_t$~: arbre de d\'ecision de profondeur max 4 (petit arbre, appel\'e \textit{weak learner})
\end{itemize}
\end{formulebox}

\subsection{Construction de chaque arbre}

\`A l'it\'eration $t$, on calcule le \textbf{gradient} de la loss par rapport aux pr\'edictions actuelles (d'o\`u le nom \textit{gradient} boosting)~:

\begin{equation}
    g_i = \frac{\partial \mathcal{L}(y_i, \hat{y}_i^{(t-1)})}{\partial \hat{y}_i^{(t-1)}}
    \qquad
    h_i = \frac{\partial^2 \mathcal{L}(y_i, \hat{y}_i^{(t-1)})}{\partial (\hat{y}_i^{(t-1)})^2}
    \label{eq:grad}
\end{equation}

Avec la log-loss (classification binaire)~:
\begin{equation}
    \mathcal{L}(y, \hat{y}) = -\left[y \cdot \ln(\hat{y}) + (1-y) \cdot \ln(1-\hat{y})\right]
\end{equation}

Le nouvel arbre $f_t$ est ajust\'e sur les r\'esidus $(g_i, h_i)$ pour minimiser~:
\begin{equation}
    \text{Gain} = \frac{1}{2}\left[\frac{G_L^2}{H_L + \lambda} + \frac{G_R^2}{H_R + \lambda} - \frac{(G_L + G_R)^2}{H_L + H_R + \lambda}\right] - \gamma
    \label{eq:gain}
\end{equation}

o\`u $G_L, G_R$ sont les sommes des gradients dans les branches gauche/droite, $\lambda$ la r\'egularisation L2, et $\gamma$ le co\^ut de complexit\'e.

\subsection{Hyperparamètres}

\begin{center}
\begin{tabular}{lcp{6cm}}
\toprule
\textbf{Param\`etre} & \textbf{Valeur} & \textbf{R\^ole} \\
\midrule
\texttt{n\_estimators} & 200 & Nombre d'arbres \\
\texttt{max\_depth} & 4 & Profondeur max (contr\^ole la complexit\'e) \\
\texttt{learning\_rate} & 0.05 & Pas d'apprentissage (petit = r\'egularisant) \\
\texttt{subsample} & 0.8 & 80\% des donn\'ees par arbre (stochastique) \\
\texttt{scale\_pos\_weight} & $\frac{N^-}{N^+}$ & Compense le d\'es\'equilibre des classes \\
\texttt{device} & \texttt{cuda} & Entra\^inement sur GPU \\
\bottomrule
\end{tabular}
\end{center}

\subsection{Calibration isotonique}

Les probabilit\'es brutes de XGBoost ne sont pas toujours bien calibr\'ees (une PD pr\'edite de 5\% ne correspond pas n\'ecessairement \`a 5\% de d\'efauts observ\'es). On applique une \textbf{calibration isotonique}~:

\begin{definitionbox}[Calibration isotonique]
On ajuste une fonction croissante par morceaux $g$ telle que~:
$$\text{PD}_{\text{calibr\'ee}} = g(\text{PD}_{\text{brute}})$$
Cette fonction est estim\'ee par r\'egression isotonique sur le jeu de validation, en minimisant~:
$$\sum_i \left(y_i - g(\hat{p}_i)\right)^2 \quad \text{sous contrainte} \quad g \text{ croissante}$$
\end{definitionbox}

\subsection{Feature importance globale}

XGBoost fournit une importance globale par variable, bas\'ee sur le \textbf{gain total} apport\'e par chaque variable dans l'ensemble des splits de tous les arbres~:

$$\text{Importance}_j = \sum_{t=1}^{T}\;\sum_{\text{split sur } j \text{ dans } f_t} \text{Gain}_{\text{split}}$$

\begin{interpretbox}[Diff\'erence avec TabNet]
L'importance XGBoost est \textbf{globale} (moyenne sur tous les clients). TabNet donne une importance \textbf{individuelle} (diff\'erente pour chaque client). Pour un audit de portefeuille, l'importance globale suffit. Pour expliquer une d\'ecision \`a un client, l'importance individuelle est n\'ecessaire.
\end{interpretbox}


% ══════════════════════════════════════════════════════════
\section{Calibration et Validation}
% ══════════════════════════════════════════════════════════

\subsection{M\'etrique principale~: AUC-ROC}

L'AUC (\textit{Area Under the ROC Curve}) mesure la capacit\'e du mod\`ele \`a \textbf{discriminer} les bons des mauvais clients~:

\begin{definitionbox}[AUC-ROC]
$$\text{AUC} = P(\hat{p}_{\text{d\'efaut}} > \hat{p}_{\text{non-d\'efaut}})$$
\begin{itemize}[noitemsep]
    \item $\text{AUC} = 0.5$~: mod\`ele al\'eatoire (aucun pouvoir discriminant)
    \item $\text{AUC} = 1.0$~: mod\`ele parfait (s\'eparation totale)
    \item $\text{AUC} \in [0.7, 0.8]$~: acceptable pour du cr\'edit corporate
    \item $\text{AUC} > 0.8$~: bon mod\`ele
\end{itemize}
\end{definitionbox}

La relation avec le coefficient de Gini~:
$$\text{Gini} = 2 \times \text{AUC} - 1$$

\subsection{Statistique de Kolmogorov-Smirnov (KS)}

Le KS mesure la \textbf{s\'eparation maximale} entre les distributions cumul\'ees des bons et des mauvais~:

$$\text{KS} = \max_t \left|F_{\text{good}}(t) - F_{\text{bad}}(t)\right|$$

Un KS $> 0.3$ est g\'en\'eralement consid\'er\'e comme satisfaisant.

\subsection{R\'esultats observ\'es}

\begin{center}
\begin{tabular}{lccc}
\toprule
\textbf{Mod\`ele} & \textbf{AUC (test)} & \textbf{Gini} & \textbf{Caract\'eristique} \\
\midrule
LR\_WoE & $\sim$0.78 & $\sim$0.56 & Transparence totale \\
TabNet & $\sim$0.81 & $\sim$0.62 & Explicabilit\'e individuelle \\
XGBoost & $\sim$0.82 & $\sim$0.64 & Performance maximale \\
\bottomrule
\end{tabular}
\end{center}

\begin{interpretbox}[Interpr\'etation des r\'esultats]
\begin{itemize}
    \item L'\'ecart \textbf{XGBoost vs TabNet est inf\'erieur \`a 1\% d'AUC}. C'est un r\'esultat classique sur donn\'ees tabulaires (Grinsztajn et al., 2022)~: le deep learning n'apporte pas de gain significatif par rapport aux m\'ethodes ensemblistes.
    \item L'\'ecart \textbf{LR vs XGBoost est de $\sim$4\% d'AUC}. C'est le prix de la transparence. La LR ne peut pas capturer les interactions et non-lin\'earit\'es que XGBoost exploite.
    \item Le triplet est \textbf{compl\'ementaire}~: la LR montre que les variables sont globalement lin\'eairement s\'eparables. XGBoost gagne 4 points en exploitant les non-lin\'earit\'es. TabNet confirme le r\'esultat de XGBoost avec un m\'ecanisme diff\'erent.
\end{itemize}
\end{interpretbox}


% ══════════════════════════════════════════════════════════
\section{Comparaison et Justification du Choix}
% ══════════════════════════════════════════════════════════

\subsection{Tableau comparatif}

\begin{center}
\begin{tabular}{p{3.5cm}ccc}
\toprule
\textbf{Crit\`ere} & \textbf{LR\_WoE} & \textbf{TabNet} & \textbf{XGBoost} \\
\midrule
AUC (test) & 0.78 & 0.81 & 0.82 \\
Interpr\'etabilit\'e globale & \textcolor{accentgreen}{\textbf{+++}} & + & ++ \\
Explicabilit\'e individuelle & + & \textcolor{accentgreen}{\textbf{+++}} & + \\
Temps d'entra\^inement (1.5M) & $\sim$10s & $\sim$3 min & $\sim$1 min \\
Nombre de param\`etres & $\sim$10 & $\sim$120\,000 & $\sim$40\,000 \\
Risque d'overfitting & Faible & Moyen & Moyen \\
Stabilit\'e temporelle & \textcolor{accentgreen}{\textbf{+++}} & + & ++ \\
Scorecard (points) & \textcolor{accentgreen}{\textbf{Oui}} & Non & Non \\
GPU requis & Non & Oui & Optionnel \\
\bottomrule
\end{tabular}
\end{center}

\subsection{Pourquoi garder les trois~?}

\begin{enumerate}
    \item \textbf{Model Risk Management (BCE/EBA)}~: le r\'egulateur exige la comparaison de plusieurs familles pour d\'etecter le \textit{model risk}. Si les trois mod\`eles s'accordent, on peut \^etre confiant. S'ils divergent, c'est un signal d'alerte.

    \item \textbf{Chaque mod\`ele a un usage diff\'erent}~:
    \begin{itemize}[noitemsep]
        \item \textbf{LR\_WoE}~: scorecard r\'eglementaire, audit, refus de cr\'edit avec justification transparente
        \item \textbf{XGBoost}~: mod\`ele de production pour le provisionnement IFRS~9 (meilleure performance)
        \item \textbf{TabNet}~: explicabilit\'e individuelle pour le comit\'e de cr\'edit (``pourquoi ce client~?'')
    \end{itemize}

    \item \textbf{Benchmark honnête}~: sans TabNet, un r\'eviseur pourrait demander ``pourquoi n'avez-vous pas test\'e le deep learning~?''. Maintenant on peut r\'epondre~: ``on a test\'e, XGBoost fait aussi bien pour 100$\times$ moins cher''.
\end{enumerate}

\subsection{Le deep learning est-il utile sur des donn\'ees tabulaires~?}

La litt\'erature r\'ecente (Grinsztajn et al., 2022 ; Borisov et al., 2022) montre de mani\`ere r\'ep\'et\'ee que~:

\begin{quote}
\textit{``Sur des donn\'ees tabulaires, les m\'ethodes ensemblistes \`a base d'arbres (XGBoost, LightGBM, CatBoost) \'egalent ou d\'epassent le deep learning dans la majorit\'e des cas.''}
\end{quote}

Les raisons principales~:
\begin{itemize}
    \item Les donn\'ees tabulaires n'ont \textbf{pas de structure spatiale ou s\'equentielle} (contrairement aux images ou au texte)
    \item Les arbres sont naturellement adapt\'es aux \textbf{interactions entre variables} et aux \textbf{seuils}
    \item Le deep learning a besoin de beaucoup plus de donn\'ees pour g\'en\'eraliser
\end{itemize}

\textbf{Notre r\'esultat confirme cette observation}~: XGBoost (AUC $\approx 0.82$) $\approx$ TabNet (AUC $\approx 0.81$) sur 1.5M lignes de cr\'edit corporate.


% ══════════════════════════════════════════════════════════
\section{Architecture Train / Serve}
% ══════════════════════════════════════════════════════════

\subsection{Phase d'entra\^inement (offline)}

\begin{center}
\begin{tikzpicture}[
    box/.style={draw=steelblue, rounded corners=4pt, minimum width=3.5cm, minimum height=1cm, fill=steelblue!8, font=\small, text=darkblue, align=center},
    data/.style={draw=accentgreen!70!black, rounded corners=4pt, minimum width=3cm, minimum height=0.8cm, fill=accentgreen!8, font=\small},
    file/.style={draw=accentorange, rounded corners=4pt, minimum width=3cm, minimum height=0.8cm, fill=accentorange!8, font=\small},
    arrow/.style={-Stealth, thick, steelblue!70}
]

\node[data] (dgp) at (0,0) {DGP v4.5\\seed=42};
\node[box] (gen) at (4,0) {G\'en\'eration\\1.5M lignes};
\node[box] (train) at (8,0) {Entra\^inement\\3 mod\`eles};
\node[file] (save) at (12,0) {pd\_suite\\.joblib};

\draw[arrow] (dgp) -- (gen);
\draw[arrow] (gen) -- (train);
\draw[arrow] (train) -- (save);

\node[below=0.2cm of gen, font=\tiny, text=gray] {$\sim$50s};
\node[below=0.2cm of train, font=\tiny, text=gray] {$\sim$5 min (GPU)};
\node[below=0.2cm of save, font=\tiny, text=gray] {$\sim$50 MB};

\end{tikzpicture}
\end{center}

Commande~:
\begin{verbatim}
python -m ifrs9_cockpit.training.train --n-clients 1500000
python -m ifrs9_cockpit.training.train --models LR_WoE,XGBoost  # partiel
\end{verbatim}

\subsection{Phase de service (dashboard)}

\begin{center}
\begin{tikzpicture}[
    box/.style={draw=steelblue, rounded corners=4pt, minimum width=3.5cm, minimum height=1cm, fill=steelblue!8, font=\small, text=darkblue, align=center},
    data/.style={draw=accentgreen!70!black, rounded corners=4pt, minimum width=3cm, minimum height=0.8cm, fill=accentgreen!8, font=\small},
    file/.style={draw=accentorange, rounded corners=4pt, minimum width=3cm, minimum height=0.8cm, fill=accentorange!8, font=\small},
    arrow/.style={-Stealth, thick, steelblue!70}
]

\node[file] (model) at (0,0) {pd\_suite\\.joblib};
\node[data] (port) at (4,0) {Portfolio\\30K (seed=123)};
\node[box] (pred) at (8,0) {Pr\'ediction\\PD};
\node[box] (dash) at (12,0) {Dashboard\\Streamlit};

\draw[arrow] (model) -- (pred);
\draw[arrow] (port) -- (pred);
\draw[arrow] (pred) -- (dash);

\node[below=0.2cm of pred, font=\tiny, text=gray] {$\sim$1s};
\node[below=0.2cm of dash, font=\tiny, text=gray] {temps r\'eel};

\end{tikzpicture}
\end{center}

Le dashboard d\'etecte automatiquement le fichier \texttt{pd\_suite.joblib}. S'il existe, les mod\`eles pr\'e-entra\^in\'es sont charg\'es ($\sim$1s). Sinon, un entra\^inement rapide sur les 30\,000 lignes du portfolio est effectu\'e ($\sim$3.5s).


% ══════════════════════════════════════════════════════════
\section{Robustesse du Pipeline}
% ══════════════════════════════════════════════════════════

\subsection{Contrat de donn\'ees}

Le pipeline impose un \textbf{contrat de donn\'ees strict} avant tout entra\^inement~:

\begin{enumerate}
    \item \textbf{Enums strictes}~: les valeurs cat\'egorielles (\texttt{sector}, \texttt{loan\_type}) sont valid\'ees contre des listes ferm\'ees. Toute valeur inconnue d\'eclenche une erreur.
    \item \textbf{Clipping des outliers}~: les variables num\'eriques sont born\'ees~:
    \begin{center}
    \begin{tabular}{lcc}
    \toprule
    \textbf{Variable} & \textbf{Min} & \textbf{Max} \\
    \midrule
    \texttt{debt\_ratio} & 0.0 & 1.5 \\
    \texttt{credit\_score} & 300 & 850 \\
    \texttt{utilization\_rate} & 0.0 & 1.2 \\
    \bottomrule
    \end{tabular}
    \end{center}
    \item \textbf{Feature engineering automatique}~: \texttt{loan\_to\_revenue} et \texttt{collateral\_coverage} sont calcul\'es \`a la vol\'ee.
    \item \textbf{Alertes NaN}~: si plus de 5\% d'une colonne est manquant, un warning est \'emis.
\end{enumerate}

\subsection{S\'eparation train / test}

\begin{itemize}[noitemsep]
    \item Split 80\% / 20\%, stratifi\'e sur la variable cible (\texttt{default\_flag})
    \item La stratification garantit le m\^eme taux de d\'efaut dans les deux ensembles
    \item Toutes les m\'etriques rapport\'ees sont sur le \textbf{jeu de test} (jamais vu pendant l'entra\^inement)
\end{itemize}


% ══════════════════════════════════════════════════════════
\section{Limites et Perspectives}
% ══════════════════════════════════════════════════════════

\subsection{Limites actuelles}

\begin{enumerate}
    \item \textbf{Donn\'ees synth\'etiques}~: le DGP, aussi sophistiqu\'e soit-il, ne remplace pas des donn\'ees r\'eelles. En production, il faudrait valider sur un historique de d\'efauts observ\'es.

    \item \textbf{Pas de validation out-of-time (OOT)}~: sur donn\'ees synth\'etiques stationnaires, l'OOT n'a pas de sens. En production, c'est indispensable (entra\^iner sur 2018--2022, tester sur 2023).

    \item \textbf{Monitoring de d\'erive (PSI/CSI)}~: non impl\'ement\'e car inutile sur donn\'ees synth\'etiques. En production, il faut surveiller la stabilit\'e de la population (PSI) et des caract\'eristiques (CSI) pour d\'etecter les d\'erives.

    \item \textbf{Contrainte de signe par \'elimination}~: malgr\'e le filtrage VIF en amont (qui stabilise les signes), l'\'elimination it\'erative des $\beta > 0$ reste un dernier recours. Une optimisation sous contrainte directe (\texttt{scipy.optimize} avec \textit{bounds}) serait plus \'el\'egante.
\end{enumerate}

\subsection{Extensions possibles}

\begin{enumerate}
    \item \textbf{PD multi-p\'eriodes}~: mod\'eliser la PD \`a 1 an, 2 ans, \ldots, N ans via des taux de hasard~:
    $$\text{PD}_{\text{lifetime}} = 1 - \prod_{t=1}^{N} (1 - \text{PD}_t) = 1 - e^{-\sum_{t=1}^N h_t}$$
    (d\'ej\`a impl\'ement\'e dans le module ECL avec $h = -\ln(1 - \text{PD})$)

    \item \textbf{Mod\`eles de survie}~: Cox proportional hazards ou DeepSurv pour mod\'eliser le \textit{time-to-default}

    \item \textbf{Ensemble des 3 mod\`eles}~: combiner les pr\'edictions par moyenne pond\'er\'ee pourrait gagner $\sim$0.5\% d'AUC suppl\'ementaire
\end{enumerate}


% ══════════════════════════════════════════════════════════
\section{R\'ef\'erences}
% ══════════════════════════════════════════════════════════

\begin{enumerate}[label={[\arabic*]}]
    \item \textbf{Siddiqi, N.} (2006). \textit{Credit Risk Scorecards: Developing and Implementing Intelligent Credit Scoring}. Wiley.

    \item \textbf{Anderson, R.} (2007). \textit{The Credit Scoring Toolkit: Theory and Practice for Retail Credit Risk Management and Decision Automation}. Oxford University Press.

    \item \textbf{Arik, S.~\"O. \& Pfister, T.} (2019). TabNet: Attentive Interpretable Tabular Learning. \textit{arXiv:1908.07442}.

    \item \textbf{Chen, T. \& Guestrin, C.} (2016). XGBoost: A Scalable Tree Boosting System. \textit{KDD 2016}.

    \item \textbf{King, G. \& Zeng, L.} (2001). Logistic Regression in Rare Events Data. \textit{Political Analysis}, 9(2), 137--163.

    \item \textbf{Grinsztajn, L., Oyallon, E. \& Varoquaux, G.} (2022). Why do tree-based models still outperform deep learning on typical tabular data? \textit{NeurIPS 2022}.

    \item \textbf{Borisov, V. et al.} (2022). Deep Neural Networks and Tabular Data: A Survey. \textit{IEEE TNNLS}.

    \item \textbf{IFRS Foundation} (2014). IFRS 9 Financial Instruments. International Accounting Standards Board.

    \item \textbf{EBA} (2017). Guidelines on PD estimation, LGD estimation and the treatment of defaulted exposures. EBA/GL/2017/16.
\end{enumerate}


\end{document}
